\begin{table}[htbp]
\centering\small
\caption{The architecture comparison.}
\label{tab:qr_v2}
\setlength{\tabcolsep}{3pt}\footnotesize
\begin{tabular}{lrrrrr}
\toprule
Arm & params & T1-arch: \texttt{logo} vs A & T2-arch: vs raw+WeChat & resc. & broke \\
\midrule
A: conv $3{\times}3$ & 19,105 & --- & $+3.5$ (part) & 1,089 & 0.5\% \\
B: hybrid & 120,353 & $+8.7$ (part) & $+10.7$ (succ) & 1,487 & 6.4\%$^{\ddagger}$ \\
C: Restormer & 52,209 & $+8.9$ (part) & $-1.2$ (nega) & 1,023 & 7.2\%$^{\ddagger}$ \\
\bottomrule
\end{tabular}
\\[4pt]
\begin{minipage}{0.94\linewidth}\footnotesize
Registered in \texttt{PREREG\_restore\_v2.md}. Percentage points of OpenCV DFR on \QrVTwoUrls{}
held-out URLs, positive when the arm is ahead. Each cell is a two-sided Wilcoxon signed-rank test on
the per-URL differences, Benjamini--Hochberg over $m=4$ (two tests, two new arms). T1-arch's bar is
$+\QrVTwoBarOne$\,pp and T2-arch's is $+\QrVTwoBarTwo$\,pp, both fixed at registration. Arm A's
T2-arch figure is the first study's T1-restore and is shown for comparison, not re-tested.
\emph{broke} is the share of renders the control arm decoded and the restored arm did not:
$^{\ddagger}$~exceeds the registered $\QrVTwoGuard\%$ guard, which is reported whatever T1-arch and
T2-arch say.
\end{minipage}
\end{table}
